We start by considering the \emph{noiseless} setting in which we are given an observation $\lmeas \in \rset^{\dimy}$ of the first $\dimy$ coordinates of the state $\state \in \rset^{\dimx}$. We denote $\topstate = (\vecidx{\state}{}{1}, \dotsc, \vecidx{\state}{}{\dimy})$ and $\bottomstate = (\vecidx{\state}{}{\dimy+1}, \dotsc, \vecidx{\state}{}{\dimx})$. In this setting, the posterior density is given by $\filter{0}{\lmeas}{\lbottomstate_0} \propto \bwmarg{0}{\lmeas, \lbottomstate_0}$.

We address the problem of drawing approximate samples from this posterior distribution  with Sequential Monte Carlo (SMC) methods. We consider the following sequence of distributions \[ 
    \filter{s}{\lmeas}{\lstate_s} \propto \int \filter{0}{\lmeas}{\lstate_0} \fwtrans{s|0}{\lstate_0}{\lstate_s} \rmd \lstate_0 \eqsp, \quad s \in [1:n-1] \eqsp, \quad \filter{n}{\lmeas}{\lstate_n} = \normpdf(\lstate_n; \mathbf{0}, \id_{\dimx}) \eqsp,
    \]
    where $\fwtrans{s|0}{\lstate_0}{\lstate_s} = \int \prod_{\ell = 1} ^s \fwtrans{\ell}{\lstate_{\ell-1}}{\lstate_\ell} \rmd \lstate_{1:s-1}$. Thus, sampling from $\filter{n}{\lmeas}{}$ and then running a SMC algorithm with distributions $\{ \filter{n - s}{\lmeas}{} \}_{s = 0} ^{n-1}$ enables approximate sampling from the posterior.
% An SMC algorithm relies on the sequence of distributions $(\filter{s}{\lmeas}{})_{s \in[0:n]}$ with $\filter{0}\lmeas{}$ as the final distribution and such that (1) $\filter{n}{\lmeas}{}$ is tractable for sampling, and (2) sampling from $\filter{s}{\lmeas}{}$ starting from $\filter{s+1}{\lmeas}{}$ is easy. The latter is achieved either by ensuring that the discrepancy between $\filter{s+1}{\lmeas}{}$ and $\filter{s}{\lmeas}{}$ is small or by designing Markov kernels $(K_s)_{s\in[0:n-1]}$ such that $\int \filter{s+1}{\lmeas}{\rmd z} K_s(\cdot | z)$ can be used as an approximation of $\filter{s}{\lmeas}{}$ for all $s \in [0:n-1]$. While the intermediate distributions $\mu_s$ are generally not available for sampling, we can use approximate samples from $\mu_{s+1}$ to obtain approximate samples from the current distribution $\mu_s$.
%  Unfortunately, these distributions are unavailable for density evaluation up to a normalizing constant and thus cannot straightforwardly be used in an SMCS scheme. 
In the next section we show how these bridge distributions can nonetheless be used to define a FK model that targets the desired posterior. 
 The resulting procedure is summarized in \Cref{alg:algopratique} and is explained throughout the next section. 

% We now consider a generic sequence of bridge distributions that builds on the marginals of the diffusion model 
% \[
% \filter{s}{\lmeas}{\rmd \statevar{s}} \propto \pot{s}{\lmeas}{\statevar{s}} \bwmarg{s}{\rmd \statevar{s}} \eqsp, \quad \text{for all $s \in [1:n]$} \eqsp,
% \]
% \eric{Pb de notation avec \eqref{eq:posterior:noisy}... Attention à avoir des notations acceptables par tous}
% where the potentials $\pot{s}{\lmeas}{}$ are positive and are yet to be designed. The core idea of our approach is to bridge consecutive marginals $\bwmarg{s+1}{}$ and $\bwmarg{s}{}$ of the diffusion model (of which the prior $\bwmarg{0}{}$  in \cref{eq:posterior:noisy,eq:posterior:noiseless} is the terminal marginal). This bridging is achieved using the transition kernel $\bwtrans{}{s}{}{}$ and with appropriately chosen potentials, we can obtain a sequence of posteriors $(\filter{s:n}\lmeas{})_{s\in[0:n]}$ that are also easy to bridge. Note that since $\bwmarg{s}{\statevar{s}} = \int \bwmarg{s+1}{\rmd \statevar{s+1}} \bwtrans{}{s}{\statevar{s+1}}{\statevar{s}}$, it is easily seen that the sequence $(\filter{s}{\lmeas}{})_s$ satisfies the following recursion;
% for all $s \in [0:n-1]$ and $B \in \sigalgebra\big( \R^{\dimx} \big)$,
% \begin{equation} 
%     \label{eq:filter:recursion}
%     \filter{s}\lmeas{B} = \frac{1}{L_{s}} \int \1_B(\statevar{s}) \filter{s+1}\lmeas{\rmd \statevar{s+1}} \bwtrans{}{s}{\statevar{s+1}}{\rmd \statevar{s}} \dubpot{s:s+1}\lmeas{\statevar{s:s+1}}\eqsp,
% \end{equation}
% where $\smash{\dubpot{s:s+1}{\lmeas}{\statevar{s:s+1}} =  \pot{s}{\lmeas}{\statevar{s}}} / \pot{s+1}{\lmeas}{\statevar{s+1}}$ and $L_{s}$ is its normalizing constant.
\begin{algorithm2e}
    \caption{\algo}
    \label{alg:algopratique} 
    \KwInput{Number of particles $N$}
    \KwOutput{$\particle{1:N}{0}$, $\weight{1:N}{0}$}
    Set $\weight{n}{1:N} = N^{-1}$. Draw $\particle{1:N}{n} \sim \normdist(0, \idm)$. \\
    \For{$s \in \intvect{n}{1}$}{
        Draw $\anc{i}{s-1} \sim \categorical\big(\{ \weight{i}{s-1} \}_{i = 1} ^N \big)$ and 
        set $\uweight{i}{s-2} = 1$.\\
        \For{$\ell \in \intvect{1}{\dimg}$}{
            \uIf{$\ell \leq \dimy$ and $\tau_\ell \leq s$}{
                Draw $\particle{i}{s - 1, \ell} \sim \bwopt{s-1|s, \ell}{\tilde{y}_\ell}{\particle{\anc{i}{s-1}}{s, \ell}}{\cdot}$ and set $\uweight{i}{s-2} = \uweight{i}{s-2} \uweight{i}{s-2, \ell}(\particle{i}{s-1, \ell})$.\\
            }
            \Else{
                Draw $\particle{i}{s - 1, \ell} = \bwtrans{}{s, \ell}{\particle{\anc{i}{s -1}}{s-1, \ell}}{\cdot}$.
            }
        }
    }
\end{algorithm2e}
\paragraph{A Feynman--Kac model based on the diffused posterior.}
We now seek natural choices of FK models by first considering the following assumption which we discuss at the end of the section.
\eric{pas une bonne idée de le mettre en hypothèse, on a l'impression que tu as besoin d'un peu moins que ça}
% \tcr{\begin{hypA} 
%     \label{assp:perfect}
%         For all $\statevar{0:n} \in (\R^{\dimx})^{n+1}$,  $\bwmarg{0}{\lstate_{0}} \fwtrans{1:n|0}{\statevar{0}}{\statevar{1:n}} = \bwmarg{0:n}{\statevar{0:n}}$.
% \end{hypA}
% }
\eric{ce n'est pas une hypothèse, c'est une approximation... n'a-t-on pas plutot besoin des marginales que l'on discute au dessus... on n'utilise jamais les lois jointes "globales"}
% In order for the SMC sampler to effectively sample from the target distribution, it is important to have a small discrepancy between $\filter{s+1}{\lmeas}{}$ and $\filter{s}{\lmeas}{}$ for all $s \in [0:n-1]$. Consider the probability measures $\filter{s}{\lmeas}{\rmd \lstate_s} \propto \int \filter{0}{\lmeas}{\rmd \lstate_0} \fwtrans{s|0}{\lstate_0}{\rmd \lstate_s}$ for $s \in [1:n]$ which correspond to the marginals of the forward diffusion applied to the posterior $\filter{0}{\lmeas}{}$. By design they bridge slowly between the posterior $\filter{0}{\lmeas}{}$ and a standard Gaussian if $n$ is large enough and 
For the sequence of diffused posteriors, we can write under 
\assp{\ref{assp:perfect}} that $\filter{s}{\lmeas}{\rmd \lstate_s} \propto  \pot{s}{\lmeas}{\lstate_s} \bwmarg{s}{\rmd \lstate_s}$ where $\pot{s}{\lmeas}{\lstate_s} = \int \bwtrans{}{0|s}{\lstate_s}{\lstate_0} \clipleb{1:m}{\lmeas}{\rmd \lstate_0}$. Following this insight we now derive in the noiseless scenario a practical recursion for the sequence $\{ \filter{s}{\lmeas}{} \}_{s = 0} ^n$ which will guide our choice of FK in the more general case when \assp{\ref{assp:perfect}} is not satisfied.

\textit{Noiseless case.} 
% Let $\filter{s}{\lmeas}{}$ denote the time $s$ marginal of the joint posterior $\filter{0:n}{\lmeas}{\lstate_{0:n}} \propto \bwmarg{0:n}{ \lstate_{0:n}} \delta_\lmeas(\ltopstate_{0})$, which is 
% \[ 
% \filter{s}{\lmeas}{\lstate_s} \propto   \int \bwmarg{0:n}{\chunk{\lstate}{0}{n}} \clipleb{0}{\lmeas}{\rmd \lstate_0} \rmd \chunk{\lstate}{1}{s-1, s+1:n} \propto \bwmarg{s}{\lstate_s} \pot{s}{\lmeas}{\lstate_s}
% \] 
% where $\smash{\pot{s}{\lmeas}{\lstate_s} = \int \bwtrans{}{0|s}{\lstate_s}{\lstate_0} \clipleb{0}{\lmeas}{\rmd \lstate_0}}$. 
Under \assp{\ref{assp:perfect}} we have that $\bwmarg{0:s+1}{} = \fwmarg{0:s+1}{}$ which implies in particular that $\bwtrans{}{0|s}{\lstate_s}{\lstate_0} \bwmarg{s}{\lstate_s} = \bwmarg{0}{\lstate_0} \fwtrans{s|0}{\lstate_0}{\lstate_s} = \int \bwmarg{0}{\lstate_0}  \revbridge{s|0}{\lstate_{s+1}, \lstate_0}{\lstate_s} \fwtrans{s+1|0}{\lstate_0}{\rmd \lstate_{s+1}}$ and thus 
\begin{align*} 
    \filter{s}{\lmeas}{\lstate_s} \propto \int \revbridge{s| 0}{\lstate_{s+1}, \lstate_0}{\lstate_s} \bwmarg{0}{\lstate_0} \fwtrans{s+1|0}{\lstate_0}{\rmd \lstate_{s+1}} \clipleb{1:\dimy}{\lmeas}{\rmd \lstate_0} \eqsp.
\end{align*}
Next, since $\topstate_{s}$ and $\bottomstate_{s}$ are independent conditionally on $(\topstate_{s+1}, \bottomstate_{s+1})$ and  depend only on $(\topstate_{s+1}, \topstate_{0})$ and $(\bottomstate_{s+1}, \bottomstate_{0})$ respectively (we write $\revbridge{s, \top}{}{}$ and $\revbridge{s, \bot}{}{}$ for the associated bridge kernels), we have that 
\begin{equation}
    \label{eq:noiseless_posterior:recursion} 
    \filter{s}{\lmeas}{\lstate_s} \propto \int \revbridge{s|0, \top}{\ltopstate_{s+1}, \lmeas}{\ltopstate_{s}} \revbridge{s|0, \bot}{\lbottomstate_{s+1}, \lbottomstate_{0}}{\lbottomstate_{s}} \bwmarg{0}{\lstate_0} \fwtrans{0|s+1}{\lstate_0}{\rmd \lstate_{s+1}} 
    \clipleb{1:m}{\lmeas}{\rmd \lstate_0} \eqsp.
\end{equation} 
Then, by replacing $\lbottomstate_{0}$ in the bridge kernel $\revbridge{s|0, \bot}{}{}$ with its prediction $\hat{\lbottomstate} _{0|s+1} (\lstate_{s+1})$ we obtain the following recursion, 
\begin{equation} 
    \label{eq:idealfilter}
    \filter{s}{\lmeas}{\rmd \lstate_s} \approx \int \filter{s+1}{\lmeas}{\rmd \lstate_{s+1}} \bwtrans{}{s, \bot}{\lstate_{s+1}}{\rmd \lbottomstate_{s}} \revbridge{s|0, \top}{\ltopstate_{s+1}, \lmeas}{\rmd \ltopstate_{s}} \eqsp,
\end{equation}
where we have used the definition  $\bwtrans{}{s}{\lstate_{s+1}}{\lbottomstate_s} = \revbridge{s|0}{\lbottomstate_{s+1}, \hat{\lbottomstate}_{0|s+1}(\lstate_{s+1})}{\lstate_s}$, see xxx.
Although this approximation is strictly valid only under the assumption of perfect measurements \assp{\ref{assp:perfect}}, it can still serve as a useful guide in the design of our FK model. Specifically, according to this recursion, the transition of $\lbottomstate_{s}$ should follow the backward diffusion while that of $\ltopstate_{s}$ should follow the bridge kernel conditioned on the measurement $\lmeas$. Mimicking this behavior while targeting the right posterior without \assp{\ref{assp:perfect}} can be done by considering the Feynman-Kac model defined as follows: 
\begin{align}
    \filter{n-1:n}{\lmeas}{\rmd \chunk{\lstate}{n-1}{n}} & \propto \bwmarg{n}{\rmd \lstate_n} \bwtrans{}{n-1, \bot}{\lstate_{n}}{\lbottomstate_{n-1}} \revbridge{n-1|0, \top}{\ltopstate_{n}, \lmeas}{\ltopstate_{n-1}} \fwtrans{n, \top}{\lmeas}{\ltopstate_n} \eqsp, \\
    \filter{0:1}{\lmeas}{\rmd \chunk{\lstate}{0}{1}} & \propto \int \filter{1:2}{\lmeas}{\rmd \chunk{\lstate}{1}{2}} \bwtrans{}{0, \bot}{\lstate_1}{\rmd \lbottomstate_0} \delta_\lmeas(\rmd \ltopstate_0) \frac{\bwtrans{}{0, \top}{\lstate_1}{\lmeas} \bwtrans{}{1, \top}{\lstate_2}{\ltopstate_1}}{\fwtrans{1,\top}{\lmeas}{\ltopstate_1} \fwtrans{2, \top}{\ltopstate_1}{\ltopstate_2}} \eqsp, \label{eq:time1joint}
\end{align}
and for all $s \in [1:n-2]$,
\begin{multline}
    \label{eq:FKrecursion}
    \filter{s:s+1}{\lmeas}{\rmd \lstate_{s:s+1}} \\
     \propto \int \filter{s+1:s+2}{\lmeas}{\rmd \lstate_{s+1:s+2}} \bwtrans{}{s, \bot}{\lstate_{s+1}}{\rmd \lbottomstate_{s}} \revbridge{s|0, \top}{\ltopstate_{s+1}, \lmeas}{\rmd \ltopstate_s} \frac{\bwtrans{}{s+1, \top}{\lstate_{s+2}}{\ltopstate_{s+1}}}{\fwtrans{s+2, \top}{\ltopstate_{s+1}}{\ltopstate_{s+2}}} \eqsp.
\end{multline}
We show in \Cref{lem:fkposterior} in the Appendix that the $\lstate_0$ marginal in \eqref{eq:time1joint} is indeed $\filter{0}{\lmeas}{}$ and hence that the introduced FK model targets the right posterior. The particle approximation $\smash{\filter{0}{N}{} = N^{-1} \sum_{i = 1}^N \delta_{\particle{i}{0}}}$, which is detailed in \Cref{appendix:sec:FK:particle_approx}, is obtained by using Algorithm~\ref{alg:algopratique} with weight function and Markov kernel 
\begin{align} 
    \uweight{}{s}(\lstate_{s+1}, \lstate_s) & = \int \fwtrans{s|0, \top}{\lmeas}{\ltopstate_s} \bwtrans{}{s}{\lstate_{s+1}}{\rmd \ltopstate_s} \big/ \fwtrans{s+1|0}{\lmeas}{\ltopstate_{s+1}} \\
    \bwopt{s}{\lmeas}{\lstate_{s+1}}{\lstate_s} & = 
\end{align}
% The particle approximation $\filter{n-1:n}{N}{}$ of $\filter{n-1:n}{\lmeas}{}$ is obtained by means of \emph{importance sampling} by sampling from the probability measure $\smash{\bwmarg{n}{\rmd \lstate_n} \bwopt{n-1}{\lmeas}{\lstate_{n-1}}{\rmd \lstate_n}}$ where $\smash{\bwopt{n-1}{\lmeas}{\lstate_{n-1}}{\lstate_n} \propto \fwtrans{n-1|0, \top}{\lmeas}{\ltopstate_{n-1}} \bwtrans{}{n-1}{\lstate_{n-1}}{\lstate_n}}$ which is available in closed form, and then weighting accordingly. 
% The particle approximation is then obtained by applying Algorithm~\ref{alg:algopratique} with the auxiliary kernel $\bwopt{s}{\lmeas}{\lstate_{s+1}}{\rmd \lstate_s} \propto \bwtrans{}{s}{\lstate_{s+1}}{\rmd \lstate_s} \fwtrans{s|0, \top}{\lmeas}{\ltopstate_s}$ which is available in closed form, and weight function $\uweight{}{s}(\lstate_{s+1}) = \int \fwtrans{s|0}{\lmeas}{\lstate_s} \bwtrans{}{s}{\lstate_{s+1}}{\rmd \lstate_s} / \fwtrans{s+1|0}{\lmeas}{\lstate_{s+1}}$. 
% Thus, we set $\pot{s}{\lmeas}{\lstate_s, \lstate_{s+1}} = \revbridge{s, \top}{\ltopstate_{s+1}, \lmeas}{\ltopstate_s}$ 
% It can be shown under \assp{\ref{assp:perfect}} that
% \begin{equation} 
%     \label{eq:smcdiff_decomp}
%     \filter{0}\lmeas{\rmd \statevar{0}} = \displaystyle \int \delta_\lmeas(\rmd \tstatevar{0}) \prod_{s = 0}^{n-1} \fwtrans{s+1,t}{\tstatevar{s}}{\rmd \tstatevar{s+1}} \filter{0:n}{}{\rmd \bstatevar{0} | \tstatevar{0:n}}  \eqsp,
% \end{equation}
% We note however that we do not need \cref{assp:perfect} for our algorithm to be consistent.

\textit{Extension to the noisy case.} 
We will now shift our focus to the posterior, which is defined by equation \eqref{eq:posterior:noisy}. \eric{pas clair ce qu'on veut dire !} It is worth noting that the likelihood function $\pot{0}\lmeas{}$ is closely connected to the forward recursion described in equation \eqref{eq:forward_def}. Specifically, if we make the assumption that there exists a $\tau_i$ for all $i \in \intvect{1}{m}$ such that $\sqrt{\alpha_{\tau_i}} \noisestd = \sqrt{1 - \alpha_{\tau_i}} s_i$, we can conclude that $\sqrt{\alpha_{\tau_i}} \lmeas_{i} / s_i = \sqrt{\alpha_{\tau_i}} \statevar{0, i} + \sqrt{1 - \alpha_{\tau_i}} \noise_i$ for all $i \in \intvect{1}{m}$. This, in turn, allows us to express the potential function $\pot{0}{\lmeas}{}$ in terms of forward kernels. More specifically, for all $i \in \intvect{1}{m}$, we can write $\normpdf(\lmeas_{i}; \statevar{0, i}, \noisestd^2) = \normpdf(\tilde{\lmeas}_{i}; \sqrt{\alpha{\tau_i}} \statevar{0, i}, (1 - \alpha_{\tau_i}))$, where $\tilde{\lmeas}_{i} = \sqrt{\alpha{\tau_i}} \lmeas_{i} / s_i$.
\begin{equation}
    \label{eq:potential:fwd_ker_decomposition}
    \pot{0}\lmeas{\lstate_{0}}  \propto \prod_{i = 1}^{m} \fwtrans{\tau_i | 0, i}{\statevar{0,i}}{\tilde{\lmeas}_{i}} \eqsp.
\end{equation}
This allows us to establish that under \assp{\ref{assp:perfect}} the posterior \eqref{eq:posterior:noisy} is the time $0$ marginal posterior of an inverse problem on the joint states $\chunk{\state}{0}{n}$ with noiseless observations $(\tilde{\lmeas}_{1}, \dotsc, \tilde{\lmeas}_{m})$ at steps $(\tau_1, \dotsc, \tau_m)$.
\begin{proposition} 
    \label{prop:noisytonoiseless}
    Assume \assp{\ref{assp:perfect}} and that $\tau_0 \eqdef 0 < \tau_1 < \dotsc < \tau_m$. Then it holds that 
    \begin{equation} 
        \noisyfilter{0}{\lmeas}{\rmd \statevar{0}} \propto \int \bwmarg{\tau_m}{\statevar{\tau_m}} \clipleb{m}{\tilde\lmeas}{\rmd \statevar{\tau_m}} \prod_{i = 0}^{m-1} \bwtrans{}{\tau_i | \tau_{i + 1}}{\statevar{\tau_{i+1}}}{\statevar{\tau_i}} \clipleb{i}{\tilde\lmeas}{\rmd \statevar{\tau_i}} \eqsp.
    \end{equation}
\end{proposition}
The proof and precise definition of the joint posterior can be found in \Cref{proof:noisytonoiseless}. 
This allows us to extend the noiseless case to adapt to \eqref{eq:posterior:noisy}. We obtain the adapted version \eqref{eq:noiseless_posterior:recursion} in Lemma~\ref{lem:decomposition_noisyposterior} which by replacing again $\lstate_0$ with its prediction $\hat{\lstate}_{0|s+1}(\lstate_{s+1})$ we obtain the following recursion for the time $s$ marginal posteriors; for all $k \in [1:m]$ (we let $\tau_{m+1} = n$) and $s \in (\tau_k, \tau_{k+1})$
\begin{equation} 
    \filter{s}{\lmeas}{\rmd \lstate_s} \approx \int \filter{s+1}{\lmeas}{\rmd \lstate_{s+1}} \left\{ \prod_{\ell = k+1} ^{\dimx} \bwtrans{}{s, \ell}{\lstate_{s+1, \ell}}{\rmd \lstate_{s, \ell}} \prod_{i = 1}^k \revbridge{s|\tau_i, i}{\lstate_{s+1, i}, \tilde\lmeas_i}{\rmd \lstate_{s, i}} \right\} \eqsp,
\end{equation}
and when $s = \tau_k$ we obtain
\begin{equation} 
    \label{eq:noisy:perfect:at_tau}
    \filter{s}{\lmeas}{\rmd \lstate_s} \approx \int \filter{s+1}{\lmeas}{\rmd \lstate_{s+1}} \left\{ \prod_{\ell = k+1} ^{\dimx} \bwtrans{}{s, \ell}{\lstate_{s+1, \ell}}{\rmd \lstate_{s, \ell}} \prod_{i = 1}^{k-1} \revbridge{s|\tau_i, i}{\lstate_{s+1, i}, \tilde\lmeas_i}{\rmd \lstate_{s, i}} \right\}  \delta_{\tilde\lmeas_k}(\rmd \lstate_{s, k})\eqsp.
\end{equation}
Therefore, proceeding in the same way as in the noiseless case, we can mimic this model. We start first by defining for all $k \in [1:m]$ and $s \in (\tau_k, \tau_{k+1})$
\begin{align}
    \bbridgeker{s}{\lstate_{s+1}, \tilde\lmeas}{\rmd \lstate_{s}} \eqdef \prod_{\ell = k+1} ^{\dimx} \bwtrans{}{s, \ell}{\lstate_{s+1}}{\rmd \lstate_{s, \ell}} \prod_{i = 1}^{k} \revbridge{s|\tau_i, i}{\lstate_{s+1, i}, \tilde\lmeas_i}{\rmd \lstate_{s, i}} \eqsp.
\end{align}
The case where $s = \tau_k$ requires special attention, since the ideal case \eqref{eq:noisy:perfect:at_tau} involves a $\delta_{\tilde\lmeas_\ell}$. To enforce this constraint while still being able to have a density w.r.t $\bwtrans{}{0:n}{}{}$, we consider a parameter $\lambda_\delta^2 > 0$ and define
$\tilde{q}_{\tau_k}(\rmd \lstate_{\tau_k, k}|\lstate_{\tau_k + 1, k}, \tilde\lmeas_k) \eqdef {\normpdf(\tilde\lmeas_k; \lstate_{\tau_k, k} ,\lambda_\delta^2)\fwtrans{\tau_k+1, k}{\lstate_{\tau_k,k}}{\lstate_{\tau_k+1,k}}} / \fwtrans{\tau_k+1|\tau_k, k}{\tilde\lmeas_k}{\lstate_{\tau_k+1,k}}$. We can finally define
\begin{multline}
    \bbridgeker{\tau_k}{\lstate_{\tau_k+1}, \tilde\lmeas}{\rmd \lstate_{\tau_k}} \eqdef \tilde{q}_{\tau_k}(\rmd \lstate_{\tau_k, k}|\lstate_{\tau_k + 1, k}, \tilde\lmeas_k)\prod_{\ell = k+1} ^{\dimx} \bwtrans{}{\tau_k, \ell}{\lstate_{\tau_k+1}}{\rmd \lstate_{\tau_k, \ell}} \\
    \prod_{i = 1}^{k-1} \revbridge{\tau_k|\tau_i, i}{\lstate_{\tau_k+1, i}, \tilde\lmeas_i}{\rmd \lstate_{\tau_k, i}} \eqsp.
\end{multline}
We are now equipped to establish the Feynman-Kac Model that mimics the ideal case. We define
\begin{equation*}
    \filter{n-1:n}{\lmeas}{\rmd \chunk{\lstate}{n-1}{n}} \propto \bwmarg{n}{\rmd \lstate_n} \bbridgeker{s}{\lstate_{n}, \tilde\lmeas}{\rmd \lstate_{n-1}} \prod_{i = 1}^\dimy \fwtrans{n, i}{\tilde\lmeas_i}{\lstate_{n, i}} \eqsp,
\end{equation*}
and for all $k \in [1:m]$ and $s \in [\tau_k, \tau_{k+1})$,
\begin{multline}
    \label{eq:FKrecursion}
    \filter{s:s+1}{\lmeas}{\rmd \lstate_{s:s+1}} 
     \propto \int \filter{s+1:s+2}{\lmeas}{\rmd \lstate_{s+1:s+2}} \bbridgeker{s}{\lstate_{s+1}, \tilde\lmeas}{\rmd \lstate_{s}}   \prod_{i=1}^{k}\frac{\bwtrans{}{s+1, i}{\lstate_{s+2}}{\lstate_{s+1,i}}}{\fwtrans{s+2, i}{\lstate_{s+1,i}}{\lstate_{s+2,i}}}  
      \eqsp.
\end{multline}
Therefore, a particle approximation is then obtained by applying Algorithm~\ref{alg:algopratique} with the choice of transition kernels and weight functions illustrated in \cref{tab:wandkerns}. All of them can be calculated in close form, as shown in \cref{appendix:explicitwandker}. Note that the algorithm coincides with the noiseless case by setting $\tau_\ell = 0$ for all $\ell \in \intvect{1}{\dimy}$.
\begin{table}[H]
    \centering
    \begin{tabular}{|c|c|c|c|}
    \hline
         & $\bwopt{s, \ell}{\tilde{\lmeas}_\ell}{\lstate_{s+1}}{\rmd \lstate_{s, \ell}}$ & $ \uweight{}{s, \ell}(\lstate_{s+1})$ \\\hline
        $s > \tau_\ell$ & $\bwtrans{}{s, \ell}{\lstate_{s+1}}{\rmd \lstate_{s, \ell}} \fwtrans{s|\tau_{\ell}, \ell}{\tilde{\lmeas}_\ell}{\lstate_{s, \ell}}$ & $\frac{\int \bwtrans{}{s}{\lstate_{s+1}}{\rmd \lstate_{s, \ell}}\fwtrans{s|\tau_{\ell}}{\tilde{\lmeas}_\ell}{\lstate_s}}{\fwtrans{s+1|\tau_{\ell}, \ell}{\lmeas}{\lstate_{s+1, \ell}}}$\\[1.5ex]\hline
        $s = \tau_\ell$ & $\bwtrans{}{s, \ell}{\lstate_{s+1}}{\rmd \lstate_{s, \ell}} \normpdf(\lstate_{\tau_\ell}; \tilde{\lmeas}_\ell, \lambda_{\delta}^2)$ & $\frac{\int \bwtrans{}{s}{\lstate_{s+1}}{\rmd \lstate_{s, \ell}}\normpdf(\lstate_{\tau_\ell}; \tilde{\lmeas}_\ell, \lambda_{\delta}^2)}{\fwtrans{s+1|\tau_{\ell}, \ell}{\lmeas}{\lstate_{s+1, \ell}}}$\\[1.5ex]\hline
        $s < \tau_\ell$& 1 & 1 \\\hline
    \end{tabular}
    \caption{Kernel and weight functions for each $\ell \in \intvect{1}{\dimy}$ and $s \in \intvect{1}{n}$. The explicit expression are given in \cref{appendix:explicitwandker}.}
    \label{tab:wandkerns}
\end{table}
% We extend \eqref{eq:noiseless:potential} for each coordinate $\ell \in \intvect{1}{\dimy}$, leading to 
% \begin{align*}
%     \pot{s, \ell}{\tilde\lmeas_{\ell}}{\lstate_s, \lstate_{s+1}} &= \fwtrans{s|\tau_{\ell}, \ell}{\tilde{\lmeas}_{\ell}}{\lstate_{s, i}} \fwtrans{s+1, \ell}{\lstate_{s, \ell}}{\lstate_{s+1, \ell}} \big/ \bwtrans{}{s,\ell}{\lstate_{s+1}}{\lstate_{s,\ell}} \quad & \text{if} \quad s > \tau_{\ell} \eqsp,\\
%     \pot{\tau_\ell, \ell}{\tilde{\lmeas}_{\tau_\ell}}{\lstate_{\tau_\ell}, \lstate_{\tau_\ell+1}} &=  \normpdf(\lstate_{\tau_\ell}; \tilde{\lmeas}_\ell, \lambda_{\delta}^2)\fwtrans{\tau_{\ell}+1, \ell}{\lstate_{\tau_{\ell}, \ell}}{\lstate_{\tau_{\ell} + 1, \ell}} \big/ \bwtrans{}{\tau_{\ell},\ell}{\lstate_{\tau_{\ell}+1}}{\lstate_{\tau_{\ell},\ell}} \quad & \text{if} \quad s = \tau_{\ell} \eqsp,
% \end{align*}
% with $\lambda_{\delta}^2 \in \rset_{>0}$ and $\pot{s, \ell}{\tilde\lmeas_{\ell}}{\lstate_s, \lstate_{s+1}} = 1$ if $s < \tau_{\ell}$. The parameter $\lambda_{\delta}^2$ should be chosen small to mimic the effect of the deltas in \cref{prop:noisytonoiseless}. If we consider then the potential $\pot{s}{\lmeas}{\lstate_s, \lstate_{s+1}} = \prod_{\ell=1}^{\dimy} \pot{s, \ell}{\tilde{\lmeas}_{\ell}}{\lstate_s, \lstate_{s+1}}$,  \eqref{eq:filter:recursion} becomes for all $k \in [1:m]$ and $s \in (\tau_k, \tau_{k+1})$, 
% \begin{align*}
%     \filter{s}{\lmeas}{\rmd \lstate_s} \propto \int \filter{s+1:s+2}{\lmeas}{\rmd (\lstate_{s+1:s+2})} &\prod_{\ell \notin \m{J}_s}\bwtrans{}{s, \ell}{\lstate_{s+1}}{\rmd \lstate_{s, \ell}} \\
%     &\prod_{\ell \in \m{J}_s} \revbridge{s|\tau_{\ell}, \ell}{\lstate_{s+1, \ell}, \tilde{\lmeas}_{\ell}}{\rmd \lstate_{s, \ell}} \frac{\bwtrans{}{s+1, \ell}{\lstate_{s+2}}{\lstate_{s+1, \ell}}}{\fwtrans{s+2, \ell}{\lstate_{s+1, \ell}}{\lstate_{s+2, \ell}}} \eqsp,
% \end{align*}
% where $\m{J}_s = \{i \in \intvect{1}{m} | \tau_{i} < s\}$. In the same way as in the noiseless case, the particle approximation is then obtained by applying Algorithm~\ref{alg:algopratique} with the auxiliary kernel $\bwopt{s, \ell}{\tilde{\lmeas}_\ell}{\lstate_{s+1}}{\rmd \lstate_{s, \ell}} \propto \bwtrans{}{s, \ell}{\lstate_{s+1}}{\rmd \lstate_{s, \ell}} \fwtrans{s|\tau_{\ell}, \ell}{\tilde{\lmeas}_\ell}{\lstate_{s, \ell}}$ and $\bwopt{\tau_{\ell}, \ell}{\tilde{\lmeas}_\ell}{\lstate_{\tau_{\ell}+1}}{\rmd \lstate_{\tau_{\ell}, \ell}} \propto \normpdf(\lstate_{\tau_\ell}; \tilde{\lmeas}_\ell, \lambda_{\delta}^2) \fwtrans{\tau_{\ell}|\tau_{\ell}, \ell}{\tilde{\lmeas}_\ell}{\lstate_{\tau_{\ell}, \ell}}$ which are available in closed form, and weight function $\uweight{}{s}(\lstate_{s+1}) = \prod_{\ell=1}^{\dimy} \uweight{}{s, \ell}(\lstate_{s+1})$ with $\uweight{}{s, \ell}(\lstate_{s+1}) = \int \fwtrans{s|\tau_{\ell}}{\tilde{\lmeas}_\ell}{\lstate_s} \bwtrans{}{s}{\lstate_{s+1}}{\rmd \lstate_{s, \ell}} / \fwtrans{s+1|\tau_{\ell}}{\lmeas}{\lstate_{s+1}}$. 
% In the particular case where $\nlessc = \emptyset$ we get that $\filter{0}\lmeas{\rmd \statevar{0}} = \int \filter{\tau_1}\lmeas{\rmd \statevar{\tau_1}} \bwtrans{}{0 | \tau_1}{\statevar{\tau_1}}{\rmd \statevar{0}}$ where $\filter{\tau_1}\lmeas{\rmd \statevar{\tau_1}} \propto \bwmarg{\tau_1}{\statevar{\tau_1}} \clipleb{\m{J}_{\tau_1}}{\lmeas}{\rmd \statevar{\tau_1}}$ and thus sampling from $\filter{0}\lmeas{}$ breaks down to sampling from the noiseless inverse problem posterior $\filter{\tau_1}\lmeas{}$ and then applying the backward diffusion.
% \begin{lemma} 
%     % \label{lem:decomposition_noisyposterior}
%     The marginals of the joint noisy posterior satisfy for all $s \in [1:n-1]$, 
%     \begin{equation} 
%         \filter{s}{\lmeas}{\lstate_s} \propto \int \bwmarg{0}{\rmd \lstate_0} \fwtrans{s+1|0}{\lstate_0}{\rmd \lstate_{s+1}} \revbridge{s,\bot}{\lbottomstate_{s+1}, \lbottomstate_{0}}{\lbottomstate_s} \prod_{\ell \in \m{J}_{\leq s}} \revbridge{s|\tau_i}{\ltopstate_{s+1,i}, \ltopstate_{0,i}}{\ltopstate_{s,i}}
%     \end{equation} 
%     where $\m{J}_{\leq s} \eqdef \{ i \in [1:n]: \tau_i \leq s \}$.
% \end{lemma}
% \textbf{old...} Given an observation $y \in \R^{\dimy}$, we aim at sampling approximately from the diffusion posterior
% \begin{equation} 
%     \label{eq:posterior}
%     \filter{0}\lmeas{\rmd x_0} =  \frac{\pot{0}\lmeas{x_0} \bwmarg{0}{\rmd x_0}}{\int \pot{0}\lmeas{\tilde{x}_0} \bwmarg{0}{\rmd \tilde{x}_0}} \eqsp.
% \end{equation}
% %which is simply the law of $X$ sampled from the diffusion model given the observation $y$ from the inverse problem model \eqref{eq:obs:model}. 
% While $\pot{0}\lmeas{}$ admits a simple form, $\bwmarg{0}{}$ does not, and as a result  the posterior \eqref{eq:posterior} is  intractable for sampling and density evaluation. 
% Furthermore, $\filter{0}\lmeas{}$ may be a high-dimensional multi-modal distribution and due to slow mixing, using standard MCMC techniques can turn out to be inefficient even if we could evaluate its unnormalized density. \emph{Sequential Monte Carlo samplers} (SMCS) \cite{del2006sequential} on the other hand offer an attractive alternative by making use of \emph{bridge distributions} that help simplify the sampling problem. In order to sample from the posterior $\filter{0}\lmeas{}$ an SMCS relies on a path of distributions $(\mu_s)_{s\in[0:n]}$ with $\mu_0 = \filter{0}\lmeas{}$ and such that (1) $\mu_n$ is tractable for sampling, (2)  sampling from $\mu_s$ starting from $\mu_{s+1}$ is easy. The latter is achieved either by making sure that the discrepancy between $\mu_{s+1}$ and $\mu_s$ is small or by designing Markov kernels $(K_s)_{s\in[0:n-1]}$ that are such that $\mu_{s+1} K_s$ is used as an approximation of $\mu_s$ for all $s \in [0:n-1]$.
% The intermediate distributions $\mu_s$ are generally not available for sampling and the idea is to use approximate samples from $\mu_{s+1}$ to obtain approximate samples from the current distribution $\mu_s$. 

% Let us sktech a simplified version of a generic SMCS and let $K_{s}$ be a $\mu_s$-invariant Markov kernel, i.e. $\mu_s K_s = \mu_s$, and  write $\mu_{s} = \tilde\mu_{s} / \int \rmd \tilde\mu_{s}$ to emphasize that the normalizing constants of the bridge distributions are not required. We then have that \[ \mu_s(\rmd x_s) = \mu_s K_s(\rmd x_s) \propto \int \frac{\rmd \tilde\mu_s}{ \rmd \tilde\mu_{s+1}}(x_{s+1}) \mu_{s+1}(\rmd x_{s+1})K_s(x_{s+1}, \rmd x_s) \eqsp,\] 
% and thus, if an empirical approximation $\mu^N _{s+1} = N^{-1} \sum_{i = 1}^N \delta_{\particle{i}{s+1}}$ of $\mu_{s+1}$ is available we can introduce the following approximation of $\mu _{s}$:
% \begin{equation}
%     \label{eq:mixture}
%     \tilde\mu^N_s(\rmd x_s) = \sum_{i = 1} ^N \weight{i}{s+1}  K_s(\particle{i}{s+1}, \rmd x_s)\eqsp, \quad \mathrm{where} \quad \weight{i}{s+1} \propto \frac{\tilde\mu_s(\particle{i}{s+1})}{\tilde\mu_{s+1}(\particle{i}{s+1})} \eqsp, 
% \end{equation}
% and $\sum_{i = 1}^N \weight{i}{s+1} = 1$. Then, to obtain an empirical approximation $\mu^N _s$ of $\mu_s$ we sample from the mixture by drawing for $\smash{\anc{1}{s+1}, \dotsc, \anc{N}{s+1} \iid \mathrm{Categorical}(\{ \weight{j}{s+1} \}_{j = 1} ^N)}$ and set 
% \[ \mu^N _s = N^{-1} \sum_{i = 1}^N \delta_{\particle{i}{s}}, \quad \mathrm{where} \quad \particle{1}{s}, \dotsc, \particle{N}{s} \sim K_s(\particle{\anc{1}{s+1}}{s+1}, \cdot) \otimes \cdots \otimes K_s(\particle{\anc{N}{s+1}}{s+1}, \cdot) \eqsp.\]
% The resampling step, which amounts to sampling from the mixture \eqref{eq:mixture}, is crucial as it prevents the degeneracy of the normalized weights and henceforth of the particle approximations. It is shown under mild assumptions that $\mu^N _s$ converges to $\mu_s$ at the usual $N^{-1/2}$ Monte Carlo rate, see \cite[Chapter 10]{douc2014nonlinear}.

% In what follows we consider a similar strategy for drawing approximate samples from $\filter{0}\lmeas{}$. Consecutive marginals $\bwmarg{s+1}{}, \bwmarg{s}{}$ of the diffusion model (of which the prior $\bwmarg{0}{}$  in \eqref{eq:posterior} is the terminal marginal) are bridged using the transition $\bwtrans{}{s}{}{}$ and thus we can build on the marginals $(\bwmarg{s}{})_{s\in[0:n]}$ to obtain a sequence of posteriors $(\filter{s:n}\lmeas{})_{s\in[0:n]}$ that are also easy to bridge. Our intermediate posteriors are
% \begin{equation} 
%     \label{eq:posterior:intermediate}
%     \filter{s:n}\lmeas{\rmd x_{s:n}} =  \frac{\pot{s:n}\lmeas{x_{s:n}} \bwmarg{s:n}{\rmd x_{s:n}}}{\int \pot{s:n}\lmeas{\tilde{x}_{s:n}} \bwmarg{s:n}{\rmd \tilde{x}_{s:n}}} \eqsp,
% \end{equation}
% where the $(\pot{s:n}\lmeas{})_{s\in[0:n]}$ are intermediate potentials introduced to define the transition kernels used to move particles from time $s+1$ to time $s$. 
% % yet to be defined and will be chosen so that it is possible to obtain an efficient transition from $\filter{s+1}\lmeas{}$ to $\filter{s}\lmeas{}$. 

% %With this in hand there is still one obstacle to overcome in order to apply the SMCS methodology we described above. 
% However, the intermediate posterior distributions have intractable unnormalized densities so that we cannot compute the weights $(\weight{i}{s+1})_{s\in[1:n]}$. By exploiting the fact that $\bwmarg{s}{ x_s} = \int \bwmarg{s+1}{\rmd x_{s+1}} \bwtrans{}{s}{x_{s+1}}{\rmd x_s}$ it is easily seen that the intermediate posterior distributions  satisfy the following recursion, for all $s \in [0:n-1]$ and $\forall B \in \sigalgebra\big( \left(\R^{\dimx}\right)^{n-s+1} \big)$,
% \begin{equation} 
%     \label{eq:filter:recursion}
%     \filter{s:n}\lmeas{B} = \frac{\int \1_B(x_{s:n}) \filter{s+1:n}\lmeas{\rmd x_{s+1:n}}\bwtrans{}{s}{x_{s+1}}{\rmd x_s} \dubpot{s:n}\lmeas{x_{s:n}} }{\int \filter{s+1}\lmeas{\rmd \tilde{x}_{s+1:n}} \bwtrans{}{s}{\tilde{x}_{s+1}}{\rmd \tilde{x}_s} \dubpot{s:n}\lmeas{\tilde{x}_{s:n}} }\eqsp,
% \end{equation}
% where $\dubpot{s:n}\lmeas{x_{s:n}}: (\R^{d_x})^{n-s+1} \ni x_{s:n} \mapsto \pot{s:n}\lmeas{x_{s:n}} / \pot{s+1:n}\lmeas{x_{s+1:n}}$. The expression \eqref{eq:filter:recursion} is reminiscent of the filtering recursion of \emph{hidden Markov models} and so it is still possible to apply the SMCS methodology. In the next two sections we proceed to detail our choice of potentials $(\pot{s}\lmeas{})_{s\in[1:n]}$ and the particle approximation of $\filter{0}\lmeas{}$.
% \subsection{Choice of potentials}
% The marginal at time $s$ of the joint posterior $\filter{0:n}{\lmeas}{}$ writes as 
% $
%     \filter{s}{\lmeas}{x_s} \propto \bwmarg{s}{x_s} \pot{s}{\lmeas}{x_s} \eqsp,
% $
% where 
% \[ 
%     \pot{s}{\lmeas}{\statevar{s}} = \int \bwtrans{}{0|s}{\statevar{s}}{\statevar{0}} \delta_\lmeas(\rmd \tstatevar{0}) \rmd \bstatevar{0} \eqsp.
% \]
% Under \assp{\ref{assp:perfect}} we can show that 
% \[ 
%     \pot{s}{\lmeas}{\statevar{s}} \propto \fwtrans{s|0, t}{\lmeas}{\tstatevar{s}} \int \bwmarg{0}{\lmeas, \bstatevar{0}} \fwtrans{s|0,b}{\bstatevar{0}}{\bstatevar{s}}
% \]
% \textbf{old...}
% Following \cite{trippe2023diffusion} first consider the assumption which implies in particular that $\datadistr = \bwmarg{0}{}$.
% % \begin{hypA} 
% % \label{assp:perfect}
% %     For all $\statevar{0:n} \in (\R^{\dimx})^{n+1}$,  $\datadistr(\statevar{0}) \fwtrans{1:n|0}{\statevar{0}}{\statevar{1:n}} = \bwmarg{0:n}{\statevar{0:n}}$.
% % \end{hypA}
% We consider first the case of noiseless inverse problems.
% For any $\statevar{} \in \R^{\dimx}$ we write $\statevar{} = (\tstatevar{}, \bstatevar{})$ where $\tstatevar{} \in \R^{\dimy}$ and $\bstatevar{} \in \R^{\dimx - \dimy}$. We assume that we  only observe the \emph{top} of the vector $\tstatevar{0} = \lmeas$ and our aim is to draw approximate samples from the posterior on $\bstatevar{0}$ given by $\noiselessfilter{}{\lmeas}{\rmd \statevar{0}} \propto \bwmarg{0}{\tstatevar{0}, \bstatevar{0}} \delta_\lmeas(\rmd \tstatevar{0}) \rmd \bstatevar{0}$. 
% Following the work of \cite{trippe2023diffusion}, it can be shown under \assp{\ref{assp:perfect}} (see ... for a proof) that
% \begin{equation} 
%     \label{eq:smcdiff_decomp}
%     \filter{0}\lmeas{\rmd \statevar{0}} = \displaystyle \int \delta_\lmeas(\rmd \tstatevar{0}) \prod_{s = 0}^{n-1} \fwtrans{s+1}{\tstatevar{s}}{\rmd \tstatevar{s+1}} \filter{0:n}{}{\rmd \bstatevar{0} | \tstatevar{0:n}}  \eqsp,
% \end{equation}
% where 
% \begin{equation*} 
%     \filter{0:n}{}{\rmd \bstatevar{0} | \tstatevar{0:n}} 
%     \propto \int \bwoptmarg{n}{t}{\rmd \bstatevar{n}} \prod_{s = 1}^{n-1} \bwopt{s}{b}{\tstatevar{s+1}, \bstatevar{s+1}}{\rmd \bstatevar{s}} \bwopt{s}{t}{\tstatevar{s+1}, \bstatevar{s+1}}{\tstatevar{s}}\eqsp.
% \end{equation*}
% \textcolor{red}{$\fwopt{s+1}{b}{\bstatevar{s+1}}{\rmd \bstatevar{s}}$, $\bwopt{s}{t}{\tstatevar{s+1}, \bstatevar{s+1}}{\tstatevar{s}}$ not defined (obvious but still).}
% This means that sampling from $\filter{0}\lmeas{}$ boils down to (1) sampling the forward process starting from $y$ up to time $n$ yielding ``diffused'' observations $\tstate{1:n}$, and then (2) sampling the states $\bstate{0:n}$ conditionally on $(\lmeas, \tstate{1:n})$. The second step requires sampling $\filter{0:n}{}{\cdot | \lmeas, \bstate{1:n}}$ which is intractable and a particle filter on the unobserved states $(\bstatevar{s})_{s \in [0:n]}$ is used in order to draw a \emph{single} approximate sample. 
% Interestingly, under this procedure it is not required to filter the states $(\tstatevar{s})_{s \in [1:n]}$ also. The particle approximation is obtained by applying Algorithm xx with the transitions $\{ \bwopt{s}{b}{\tstate{s+1}, \cdot}{\cdot} \}_{s \in [0:n-1]}$ and potentials $\{ \bwopt{s}{t}{\tstate{s+1}, \cdot}{\tstate{s}}\} _{s \in [0:n-1]}$. It is worth noting that the integral form \eqref{eq:smcdiff_decomp} does not hold without \assp{\ref{assp:perfect}} and as such, filtering the whole states $(\statevar{s})_{s \in [0:n]}$, instead of the unobserved states only, is necessary for obtaining a consistent empirical approximation of $\filter{0}{\lmeas}{}$. 

% For our choice of potentials we draw inspiration from smcdiff \textcolor{red}{Additional notations to define here}
% \[ 
%    \fwopt{1}{t}\lmeas{\tstatevar{1}} \prod_{s = 1}^{n-1} \fwopt{s+1}{t}{\tstatevar{s}}{\tstatevar{s+1}} = \fwopt{n|0}{t}\lmeas{\tstatevar{n}} \prod_{s = 1}^{n-1} \revbridge{s}{y,t}{\tstatevar{s+1}}{\tstatevar{s}} \eqsp,
%     \]
% and considering the intermediate posteriors initialized with $ \filter{n}{\lmeas}{\rmd \statevar{n}} \propto \fwtrans{n|0}\lmeas{\tstatevar{n}} \bwmarg{n}{\rmd \statevar{n}} $,
% \begin{equation}
%     \filter{s:n}\lmeas{\rmd \statevar{s:n}} \propto  \revbridge{s}\lmeas{\tstatevar{s+1}}{\tstatevar{s}} \bwmarg{s:n}{\rmd \statevar{s:n}} \eqsp, \quad \forall s \in [1:n-1] \eqsp,
% \end{equation}
% with $\filter{0:n}\lmeas{\rmd \statevar{0}} \propto \bwopt{0}{t}{\tstatevar{1}}\lmeas \bwmarg{0:n}{\rmd \statevar{0:n}}$ as terminal posterior. The recursion \eqref{eq:filter:recursion} can then be written as, 
% \begin{equation} 
%     \filter{s:n}{\lmeas}{\rmd \statevar{s:n}} \propto \int \filter{s+1:n}{\lmeas}{\rmd \statevar{s+1:n}} \bwopt{s}{b}{\statevar{s+1}}{\rmd \bstatevar{s}} \revbridge{s}{\lmeas, t}{\tstatevar{s+1}}{\rmd \tstatevar{s}} \frac{\bwopt{s}{t}{\statevar{s+1}}{\tstatevar{s}}}{\revbridge{s+1}{\lmeas, t}{\tstatevar{s+2}}{ \tstatevar{s+1}}} \eqsp.
% \end{equation}
% Then, running Algorithm with the transitions $(\statevar{s+1}, A) \mapsto \int \1_A(\statevar{s})\bwopt{s}{b}{\statevar{s+1}}{\rmd \bstatevar{s}}  \revbridge{s}{\lmeas, t}{\tstatevar{s+1}}{\rmd \tstatevar{s}}$ and weights $\bwopt{s}{t}{\statevar{s+1}}{\cdot} / \revbridge{s+1}{\lmeas, t}{\tstatevar{s+2}}{\cdot}$ yields a consistent approximation
% \begin{align*}
%     \filter{0}\lmeas{x_0} &  \propto \int \bwoptmarg{n}{t}{\bstatevar{n}} \fwtrans{n|0}\lmeas{ x^b _n}  \left[\prod_{s = 1}^{n-1} \revbridge{s}\lmeas{\bstatevar{s+1}}{\bstatevar{s}}
%     \bwopt{s}{}{(\bstatevar{s+1}, \tstatevar{s+1})}{(x_s^t, x^b_s)} \right]\\
%     &\times 
%     \bwopt{0}{}{(\bstatevar{1}, \tstatevar{1})}{(x^t_0, y)} \rmd x^t_{1:n} \rmd x^b_{1:n} \eqsp.
% \end{align*}
% Therefore, we recover the particle approximation of \cite{trippe2023diffusion} by starting at  and considering for $s \in \intvect{1}{n-1}$  
% \gabriel{On ne devrait pas avoir tous les q ici?}
% and \yazid{todo:il reste pas mal de choses à dire ici.}
% \begin{align*} 
%     \filter{0}\lmeas{\rmd \statevar{0:n}} & \propto \bwmarg{0:n}{\rmd \statevar{0:n}} \bwopt{0}{b}{\bstatevar{1}}{\bstatevar{0}} \delta_\lmeas(\rmd \bstatevar{0}) \eqsp.
% \end{align*}

% We now turn to the case of noisy inverse problems. As explained in \cref{sec:inverse_problem_formulation} we can consider without loss of generality a potential of the form $\pot{0}\lmeas{\tstatevar{0}} = \normpdf(y; \tstatevar{0}, \Sigma_\lmeas)$ and we remind the reader that $\Sigma_\lmeas$ is an $\dimy$-diagonal matrix with non-negative entries. To simplify the exposition we assume that all the diagonal elements are different, i.e. the diagonal entries of $\Sigma_\lmeas$ are $\eigsigy{1}^2 < \dotsc < \eigsigy{\dimy}^2$. The extension to a block of coordinates with equal variance is provided Appendix xxx. % We denote by $o_i$ the number of occurences of the singular value $\eigSigmay_i$ which satisfy $\sum_{i = 1}^m o_i = \dimy$ and 
% We let $\tau_1 < \dotsc < \tau_{m}$ be the time steps such that $\alpha_{\tau_s} \eigsigy{s}^2 = 1 - \alpha_{\tau_s}$ for $s \in \intvect{0}{m}$.
% Then, letting $\tilde{y}_i = \sqrt{\alpha _{\tau_i}} y_i$, we can write that \yazid{todo: further comments here}
% \begin{equation}
%     \label{eq:potential:fwd_ker_decomposition}
%     \pot{0}\lmeas{\tstatevar{0}} = \prod_{i = 0}^{\dimy} \normpdf(y_i; \tstatevar{0,i}, \eigsigy{i}^2) \propto 
%     \prod_{i = 0}^{\dimy} \normpdf(\tilde{y}_i; \sqrt{\alpha} _{\tau_i} \tstatevar{0,i}, 1 - \alpha_{\tau_i})\eqsp.
% \end{equation}
% %where $y^{i}$ are the coordinates of $y$ with variance $\eigSigmay_i$ and $\bstatevar{0,[i]}$ is the associated mean. 
% As a consequence,  the potential $\pot{0}\lmeas{\tstatevar{0}}$ can be written in terms of forward kernels, \[
%      \pot{0}\lmeas{\tstatevar{0}} \propto \prod_{i = 1}^{\dimy} \fwopt{\tau_i | 0}{i}{\tstatevar{0,i}}{\tilde{y}_i} \eqsp.
%      \]
% Following this decomposition we now show that under \assp{\ref{assp:perfect}} the problem of sampling from the posterior of a noisy inverse problem breaks down to that of sampling from the posteriors of some particular noiseless inverse problems.
% Let $\noiselessfilter{\tau_\dimy}{\tilde{y}_m}{\rmd \statevar{\tau_\dimy}} \propto \bwmarg{\tau_\dimy}{\statevar{\tau_\dimy}} \clipleb{\dimy}{\tilde{y}_{\dimy}}{\rmd \statevar{\tau_\dimy}}$ and  define recursively for all $s \in [1:\dimy-1]$
% \[ 
%     \noisyfilter{\tau_{s}}{\lmeas}{\rmd \statevar{\tau_s}} \propto \int \filter{\tau_{s+1}}\lmeas{\rmd \statevar{\tau_{s+1}}} \bwtrans{}{\tau_{s} | \tau_{s+1}}{\statevar{\tau_{s+1}}}{\statevar{\tau_s}} \clipleb{s}{\tilde{y}_s}{\rmd \statevar{\tau_s}} \eqsp.
% \]  
% \begin{proposition} 
%     \label{prop:noisytonoiseless}
%     Assume \assp{\ref{assp:perfect}}. Then it holds that 
%     \begin{equation} 
%         \noisyfilter{0}{\lmeas}{\rmd \statevar{0}} \propto \int \noisyfilter{\tau_1}{\lmeas}{\rmd \statevar{\tau_1}} \bwtrans{}{0|\tau_1}{\statevar{\tau_1}}{\rmd \statevar{0}} \eqsp.
%         %  \\= \int \bwmarg{\tau_{\dimy}}{\rmd \statevar{\tau_{\dimy}}} \delta_{\tilde{y}_{\dimy}}(\rmd \tstatevar{\tau_{\dimy}, \dimy})
%         %   \prod_{j = 1}^{\dimy-1} \bwtrans{}{\tau_{j} | \tau_{j+1}}{\statevar{\tau_{j+1}}}{\rmd \statevar{\tau_j}} \delta_{\tilde{y}_{j}}(\rmd \tstatevar{\tau_j, j}) \bwtrans{}{0|\tau_1}{x_{\tau_1}}{x_0} \eqsp.
%     \end{equation}
% \end{proposition}
% Following this insight, our intermediate posteriors are defined for all $i \in [1:m]$,
% \begin{align*} 
%     \filter{s:n}{\lmeas}{\rmd \statevar{s:n}} & \propto \revbridge{s}{\lmeas, i}{x_{s+1, i}}{x_{s, i}} \bwmarg{s:n}{\rmd \statevar{s:n}}, \quad s \in (\tau_i: \tau_{i+1}] \eqsp, \\
%     \filter{\tau_{i}: n}{\lmeas}{\rmd \statevar{\tau_{i}:n}} & \propto \bwopt{\tau_i}{i}{x_{\tau_i, i}}{\tilde{y}_i} \bwmarg{\tau_{i}:n}{\statevar{\tau_i : n}} \clipleb{i}{\tilde{y}_{i}}{\rmd \statevar{\tau_i : n}} \\
%     \filter{0:\tau_1}{\lmeas}{\rmd \statevar{0:\tau_1}} & \propto \pot{0}{\lmeas}{\tstatevar{0}} \bwmarg{0:\tau_1}{\rmd \statevar{0:\tau_1}}
% \end{align*}

% \begin{algorithm2e}
%     \caption{\ouralgorithm}
%     \label{alg:APF} 
%     \KwInput{$N$}
%     \KwOutput{$\particle{1:N}{0}$, $\weight{1:N}{0}$}
%     Compute time steps $\tau_{0:m}$.\\
%     Draw $\particle{1:N}{n} \sim \ndist{0}{\idm}$.\\
%     \For{$s \in [n:1]$}{
%         set $J = \{\ell | \tau_\ell < s\}$ \\
%         set $\weight^{1:N} = \uweight(\particle{1:N}{s})$
%         sample $\ancestors{1:N}{s} \sim $\\
       
%     }
%     $\anc{1:N}{\tau_1} \iid \categorical\big( \{ \weight{j}{\tau_1} \}_{j = 1}^N\big)$;\\
%     $\particle{1:N}{0} \iid \bigotimes_{j = 1}^N \bwtrans{}{0 | \tau_i}{\particle{A^j _{\tau_1}}{\tau_1}}{\cdot}$;\\
%     Compute $\uweight{i}{0} = \pot{0}{\lmeas}{\particle{i}{0}}$ and set $\weight{i}{0} = \uweight{i}{0} / \sum_{j = 1}^N \uweight{j}{0}$.
%  \end{algorithm2e}
% and
% \begin{align*} 
%     L_{1:n}(y, \rmd \tstatevar{1:n}) & = \int \bwoptmarg{n}{t}{\rmd \tilde{x}^t _n} \prod_{s = 0}^{n-1} \bwopt{s}{t}{(\tilde{x}^t _{s+1}, \tstatevar{s+1})}{\rmd \tilde{x}^t _s} \bwopt{s}{b}{(\tilde{x}^t _{s+1}, \tstatevar{s+1})}{x^b _s} \delta_\lmeas(\rmd x^b _0) \eqsp, \\
%     M^b _{1:n|0}(y, \rmd \tstatevar{1:n}) & = \frac{L_{1:n}(y, \rmd \tstatevar{1:n})}{\int \bwmarg{0}{\tilde{x} _0} \delta_\lmeas(\rmd x^b _0) \rmd \tilde{x}^t _0} \eqsp.
% \end{align*}
\paragraph{Consistency}
% \tcr{To be rewritten}
% The approximation mechanism we have just described is known as the fully adapted \emph{auxiliary particle filter} introduced in \cite{pitt1999filtering}. Crucially, we have assumed that it is possible to sample from $\bwopt{s}\lmeas{}{}$ for all $s \in [0:n-1]$ and $\filter{n}\lmeas{}$ and we now detail how this can be made possible. 
% \[ 
%     \pot{s}\lmeas{x_s}  = \normpdf(y; c_s A x_s, \Sigma_s) \eqsp,
% \]
% where $c_s \in \R$ and $\Sigma_s$ is a positive definite matrix. 
% \begin{hypA}
%     \label{assp:bounded_score}
%     For all $s \in [0:n]$ there exists $C_s < \infty$ such that $\| \epsprop{\cdot, s}{\param^\star} \|_\infty \leq C_s$. 
% \end{hypA}

% \begin{lemma}
%     \label{lem:potentials}
%     Assume \assp{\ref{assp:bounded_score}}. Consider the reverse transitions and potentials defined for any $s \in [1:n]$ by
%     \begin{align*} 
%         \bwtrans{}{s}{x_{s}, \rmd x_{s-1}}{x_{s-1}} & = \normpdf(x_{s-1}; \gamma_{1,s} x_{s} + \gamma_{2,s} \epsprop{s, x_s}{\param^\star}, \sigma_s \mathbf{I}_{\dimx}) \rmd x_{s-1}\\
%         \pot{s}\lmeas{x_s} 
%         & = \normpdf(y_s; c_s A x_s, \Sigma_\lmeas + (c_s ^2 - 1) A A^\intercal)  \eqsp,
%     \end{align*}
%     where the sequence $(c_s)_{s \in [0:n]}$ with $c_0 = 1$ satisfies 
%     \begin{equation}
%         \label{eq:coeff:assumptions}
%     \frac{c^2 _{s-1} \gamma_{1,s} ^2}{\eigSigmay_i + (c^2 _{s-1} - 1 + \sigma^2 _s) \eigA _i} \geq \frac{c^2 _s}{\eigSigmay_i + (c^2 _s - 1) \eigA _i}, \quad i \in [1:n] \eqsp.
%     \end{equation}
%     Then it holds that $\sup_{x_{s+1} \in \R^{\dimx}} \uweight{}{s+1}(x_{s+1}) < \infty$.
% \eric{c'est quoi $\sigma_i^A$ ?}
% \end{lemma}
% \begin{example}[DDPM]
%     For all $s \in [1:n]$ let $c_s = \prod_{\ell = 0}^s (1 - \beta_\ell)^{-1}$ and set $c_0 = 1$. Since for DDPM $\gamma_{1, s} = 1 / \sqrt{1 - \beta_s}$ and $\sigma_s = \sqrt{\beta_s}$, we have that $c^2 _{s-1} \gamma_{1,s} ^2 = c^2 _s$ and for all $i \in [1:\dimy]$
%     \[ 
%       \eigSigmay _i + (1 - \beta_s)(c^2 _s - 1) \eigA_i  = \eigSigmay _i + (c^2 _{s-1} - 1 + \sigma^2 _s) \eigA_i \leq \eigSigmay_i + (c^2 _s - 1) \eigA_i \eqsp,
%     \]
%     since $\beta_s \in [0,1]$, and thus \eqref{eq:coeff:assumptions} is satisfied.
% \end{example}
\textbf{do not read, not finished...} We now provide a bound on the backward KL divergence between $\filter{0}\lmeas{}$ and its particle-based estimation $\bfilter{0}{N}{}$ defined for any $B \in \sigalgebra(\R^d)$ by
\[
    \bfilter{0}{N}{B} = \pE \big[ \filter{0}{N}{B} \big]\,,
    \]
where the expectation is taken with respect to the law of the generated particle cloud $\lawSMC{0:n}$:
\begin{equation} 
    \lawSMC{0:n}(\rmd \iparticle{1:N}{0:n}, \rmd \ianc{1:N}{1:n}) = \bigg\{\prod_{i = 1}^N \bwoptmarg{n}\lmeas{\rmd \iparticle{i}{n}}{} \bigg\}\prod_{s = 1}^n \bigg\{ \prod_{k = 1}^N   \sum_{\ell = 1}^N \weight{\ell}{s} \delta_\ell (\rmd \ianc{k}{s})  \bwopt{s}\lmeas{\iparticle{\ianc{k}{s}}{s}}{\rmd \iparticle{k}{s-1}} \bigg\} \eqsp.
\end{equation}
Define the \emph{joint smoothing} distribution 
\begin{equation} 
    \filter{0:n}\lmeas{\rmd x_{0:n}} = \frac{\bwmarg{0:n}{\rmd x_{0:n}}{}  g_0(x_0)}{\int \bwmarg{0}{\rmd \tilde{x}_{0}}{}  g_0(\tilde{x}_0)} \eqsp.
\end{equation}
For all $s \in [0:n]$ we write $\normconst_s = \int \bwmarg{s}{\rmd x_s} \pot{s}\lmeas{x_s}$.
\begin{hypA} 
    \label{assp:bounded_weights}
    For all $\ell \in [1:n]$, there exists $C_\ell < \infty$ such that $\sup_{x_\ell \in \R^{\dimx}} \uweight{}{\ell-1|\ell}(x_\ell) \leq C_\ell.$
\end{hypA}

\begin{proposition} \label{lem:kldiv:filtering}
    Assume \assp{\ref{assp:bounded_weights}}. Then it holds that
    \begin{multline*}
        \kldivergence{\filter{0}\lmeas{}}{\bfilter{0}{N}{}} \leq n \log \left(1-\frac{1}{N}\right) \\
         + \frac{1}{N-1}  \sum_{s = 1}^n \int \frac{\normconst_s / \normconst_0}{\pot{s}\lmeas{z_s}} \int \bwtrans{}{0|s}{z_s}{\rmd x_0} \pot{0}\lmeas{x_0} \filter{0:n}\lmeas{\rmd z_{0:n}} + \bigo(N^{-2}) \eqsp.
    \end{multline*}
\end{proposition}
\textcolor{red}{We know the dependency on $n$ of the second term (sum + $\phi_{0:n}^y$) ? }
\begin{proof}  
    The proof is given in \Cref{proof:kldiv:filtering}.
\end{proof}

% \subsection{Application to the noiseless case.}
% In this section we consider target posteriors of the type
% \begin{equation} 
%     \filter{0}\lmeas{\rmd \bstatevar0} = \frac{\bwmarg{0}{\bstatevar0, y}  \lambda(\rmd \bstatevar0)}{\int \bwmarg{0}{\tilde{x}^t _0, y} \lambda(\rmd \tilde{x}^t _0)}
% \end{equation}
% We have that
% \begin{align*} 
%     \filter{0}\lmeas{\rmd \bstatevar0} = \int \filter{0:n}\lmeas{\tstatevar{1:n}, \rmd \bstatevar{0:n}} M^b _{1:n | 0}(y, \rmd \tstatevar{1:n}) \eqsp,
% \end{align*}
% where 
% \begin{multline*} 
%      \filter{0:n}\lmeas{x_{1:n} ^b, \rmd \bstatevar{0:n}} \\
%      = \frac{1}{L_{1:n}(y, \tstatevar{1:n})} \bwoptmarg{n}{t}{\rmd \bstatevar{n}} \prod_{s = 0}^{n-1} \bwopt{s}{t}{(\bstatevar{s+1}, \tstatevar{s+1})}{\rmd \bstatevar{s}} \bwopt{s}{b}{(\bstatevar{s+1}, \tstatevar{s+1})}{x^b _s} \delta_\lmeas(\rmd x^b _0) \eqsp,
% \end{multline*}
% and 
% \begin{align*} 
%      M^b _{1:n|0}(y, \rmd \tstatevar{1:n}) & = \frac{L_{1:n}(y, \rmd \tstatevar{1:n})}{\int \bwmarg{0}{\tilde{x}^t _0, y} \lambda(\rmd \tilde{x}^t _0)} \eqsp,\\
%      L_{1:n}(y, \rmd \tstatevar{1:n}) & = \int \bwoptmarg{n}{t}{\rmd \tilde{x}^t _n} \prod_{s = 0}^{n-1} \bwopt{s}{t}{(\tilde{x}^t _{s+1}, \tstatevar{s+1})}{\rmd \tilde{x}^t _s} \bwopt{s}{b}{(\tilde{x}^t _{s+1}, \tstatevar{s+1})}{x^b _s} \delta_\lmeas(\rmd x^b _0) \eqsp.
% \end{align*}
% Let us describe the SMCDIFF methodology developed in \cite{trippe2023diffusion}. The authors assume that the joint backward and forward processes match, i.e. $\bwmarg{0:n}{\rmd x_{0:n}} = \fwmarg{0:n}{\rmd x_{0:n}}$ since under this assumption we can write that
% \begin{align*} 
% L_{1:n}(y, \rmd \tstatevar{1:n}) = \int \fwmarg{0:n}{\rmd x_{0:n}} \delta_\lmeas(\rmd x^b _0) = \int \bwmarg{0}{\bstatevar0, y} \lambda(\rmd \bstatevar0) \int \prod_{s = 0}^{n-1} \fwtrans{s+1}{x^b _s}{\rmd \tstatevar{s+1}}
% \end{align*}
% and thus, the original posterior simplifies to
% \begin{align*} 
%     \filter{0}\lmeas{\rmd \bstatevar0} = \int \delta_\lmeas(\rmd x^b _0) \prod_{s = 0}^{n-1} \fwtrans{s+1}{x^b _s}{\rmd \tstatevar{s+1}} \filter{0:n}\lmeas{\tstatevar{1:n}, \rmd \bstatevar{0:n}} \eqsp.
% \end{align*}
% Then, in order to obtain an empirical approximation, they first sample \emph{one forward trajectory} $(\particle{b}{1}, \dotsc, \particle{b}{n}) \sim \delta_\lmeas(\rmd x^b _0) \prod_{s = 0}^{n-1} \fwtrans{s+1}{x^b _s}{\rmd \tstatevar{s+1}}$ and then approximate the conditional distribution $\filter{0:n}\lmeas{\particle{b}{1:n}, \rmd \bstatevar{0:n}}$ using particle filters. We now show that the the same approximation can be achieved without the requirement that the backward and forward match perfectly.

% We define for all $s \in [1:n]$
% \begin{align*} 
%     \pred{s}{}{\tstatevar{s-1:n}, \rmd \bstatevar{s}} & = \int \filter{s+1}\lmeas{\tstatevar{s-1:n}, \rmd \bstatevar{s+1}}  \bwtrans{}{s}{(\bstatevar{s+1}, \tstatevar{s+1})}{\rmd \bstatevars} \\
%     \filter{s}\lmeas{\tstatevar{s-1:n}, \rmd \bstatevar{s+1}} & \propto \pred{s}{}{\tstatevar{s:n}, \rmd \bstatevar{s+1}} \frac{\bwtrans{}{s}{(\bstatevar{s+1}, \tstatevar{s+1})}{x^b _s}}{\fwtrans{s}{\tstatevar{s-1}}{x^b _s}} \eqsp.
% \end{align*}
% For $s \in []$, 
% \begin{align*} 
%     \filter{n}{}{\tstatevar{n-2:n}, \rmd \bstatevarn} & = \bwmarg{n}{\rmd \bstatevarn} \frac{\bwmarg{n}{\tstatevar{n}} \bwtrans{}{n-1}{(\bstatevarn, x^b _n)}{\tstatevar{n-1}}}{\fwtrans{n}{\tstatevar{n-1}}{x^b _n} \fwtrans{n-1}{\tstatevar{n-2}}{\tstatevar{n-1}}} \\
%     \filter{0}{}{\tstatevar{0:n}, \rmd \bstatevar{1}} & \propto \pred{0}{}{\tstatevar{0:n}, \rmd \bstatevar1} \bwtrans{}{0}{(\bstatevar1, x^b _1)}{x^b _0}
% \end{align*}
% Thus, we can write that 
% \[ 
%     \filter{0}\lmeas{\rmd \bstatevar0} = \int \frac{ \delta_\lmeas(\rmd x^b _0) M^b _{1:n}(x^b _0, \rmd \tstatevar{1:n}) \pred{0}{}{\tstatevar{0:n}, \mathbf{1}} \filter{0}{}{\tstatevar{0:n}, \rmd \bstatevar{0}}}{\int\delta_\lmeas(\rmd \tilde{x}^b _0) M^b _{1:n}(\tilde{x}^b _0, \rmd \tilde{x}^b _{1:n})  \pred{0}{}{\tilde{x}^b _{0:n}, \mathbf{1}}}
% \]
% where 
% \[ 
% M^b _{1:n}(x^b _0, \rmd \tstatevar{1:n}) = \int \prod_{s = 0}^{n-1} \fwtrans{s+1}{x^b _s}{\rmd \tstatevar{s+1}}
% \]



\paragraph{Connection to SMC diff} Following the work of \cite{trippe2023diffusion}, it can be shown under \assp{\ref{assp:perfect}} that
\begin{equation} 
    \label{eq:smcdiff_decomp}
    \filter{0}\lmeas{\rmd \statevar{0}} = \displaystyle \int \delta_\lmeas(\rmd \tstatevar{0}) \prod_{s = 0}^{n-1} \fwtrans{s+1,t}{\tstatevar{s}}{\rmd \tstatevar{s+1}} \filter{0:n}{}{\rmd \bstatevar{0} | \tstatevar{0:n}}  \eqsp,
\end{equation}
where 
$
    \filter{0:n}{}{\rmd \bstatevar{0} | \tstatevar{0:n}} 
    \propto \int \bwoptmarg{n}{t}{\rmd \bstatevar{n}} \prod_{s = 1}^{n-1} \bwopt{s}{b}{\tstatevar{s+1}, \bstatevar{s+1}}{\rmd \bstatevar{s}} \bwopt{s}{t}{\tstatevar{s+1}, \bstatevar{s+1}}{\tstatevar{s}}\eqsp.
$
\textcolor{red}{$\fwopt{s+1}{b}{\bstatevar{s+1}}{\rmd \bstatevar{s}}$, $\bwopt{s}{t}{\tstatevar{s+1}, \bstatevar{s+1}}{\tstatevar{s}}$ not defined (obvious but still).}
This means that sampling from $\filter{0}\lmeas{}$ boils down to (1) sampling the forward process starting from $y$ at time $0$ and up to time $n$ yielding ``diffused'' observations $\tstate{1:n}$, and then (2) sampling the states $\bstate{0:n}$ conditionally on $(\lmeas, \tstate{1:n})$. The second step requires sampling $\filter{0:n}{}{\cdot | \lmeas, \bstate{1:n}}$ which is intractable but since it has the structure of a Feynman--Kac model with transitions $\{ \bwopt{s}{b}{\tstate{s+1}, \cdot}{\cdot} \}_{s \in [0:n-1]}$ and potentials $\{ \bwopt{s}{t}{\tstate{s+1}, \cdot}{\tstate{s}}\} _{s \in [0:n-1]}$, a particle filter can be used to obtain an empirical approximation from which a \emph{single} approximate sample is then drawn. Since \assp{\ref{assp:perfect}} is an ideal assumption that does not hold in practice, it is natural to wonder if this procedure is still theoretically grounded in realistic scenarios. With no assumption we can write that 
\begin{multline}
    \filter{0}{\lmeas}{\rmd \statevar{0}} = \int \delta_\lmeas(\rmd \tstatevar{0}) \prod_{s = 0}^{n-1} \fwtrans{s+1, t}{\tstatevar{s}}{\rmd \tstatevar{s+1}} \\
    \times \frac{1}{L_{0:n}(\lmeas)} \frac{\bwmarg{n,b}{\rmd  \bstatevar{n}} \bwmarg{n,t}{\tstatevar{n}} \prod_{s = 0}^{n-1} \bwtrans{}{s, b}{\statevar{s+1}}{\rmd \bstatevar{s}} \bwtrans{}{s, t}{\statevar{s+1}}{\tstatevar{s}}}{\prod_{s = 0}^{n-1} \fwtrans{s+1, t}{\tstatevar{s}}{\tstatevar{s+1}}}\eqsp,
\end{multline}
where $L_{0:n}(\lmeas) \eqdef \int \bwmarg{0}{\statevar{0}} \delta_\lmeas(\rmd \tstatevar{0}) \rmd \bstatevar{0}$ and we see that the remaining measure on the unobserved states $\bstatevar{0:n}$ does not integrate to one.